%% kapa: exam=1 round=21 year=2010 period=2 form=A subject=appraisal_law subject_ko=부동산관계법규 seq=11 number=51
%% source: Q-Net 감정평가사 기출문제 (cst00309) artlSeq=5024181 file=extracted/1008/5024181/2010년 제21회 감정평가사 1차시험 기출문제/2교시.hwp
%% conversion: kordoc 4.24.51 -> tools/gen1.py (auto); review=proofread-gpt-sol; flags=none
\begin{question}{51}
국토의 계획 및 이용에 관한 법령상 용도지역 안에서의 행위제한에 관한 설명으로 옳지 않은 것은?
\begin{choices}
  \item 용도지역에서의 건축물이나 그 밖의 시설의 용도ㆍ종류 및 규모 등의 제한에 관한 사항은 대통령령으로 정한다.
  \item 「산업입지 및 개발에 관한 법률」에 따른 일반산업단지에서는「산업입지 및 개발에 관한 법률」에서 정하는 바에 따른다.
  \item 농림지역 중 농업진흥지역인 경우에는「농지법」에서 정하는 바에 따른다.
  \item 자연환경보전지역 중「수도법」에 따른 상수원보호구역인 경우에는「수도법」에서 정하는 바에 따른다.
  \item 자연환경보전지역 중 수산자원보호구역인 경우에는「수산자원관리법」에서 정하는 바에 따른다.
\end{choices}
\end{question}
\answer{2}
